% =====================================================================
%  CARNET D'ENTRAÎNEMENT — Fiche 12 (Chimie organique)
%  THÈME 1 — Hydrocarbures, formules brutes et squelette carboné
%  NIVEAU DIFFICILE — ÉNONCÉS
% =====================================================================
\documentclass[11pt,a4paper]{article}
\usepackage[utf8]{inputenc}
\usepackage[T1]{fontenc}
\usepackage{lmodern}
\usepackage[french]{babel}
\usepackage{amsmath,amssymb,mathtools}
\providecommand{\degree}{^\circ}
\usepackage{icomma}
\usepackage{siunitx}
\sisetup{output-decimal-marker={,},per-mode=symbol}
\usepackage[version=4]{mhchem}
\usepackage{chemfig}
\setchemfig{atom sep=2em,bond length=0.9em,cram width=2.5pt}
\usepackage{xcolor}
\usepackage{array}
\usepackage{booktabs}
\usepackage{tabularx}
\usepackage[most]{tcolorbox}
\usepackage{enumitem}
\usepackage{tikz}
\usetikzlibrary{arrows.meta,positioning,calc}
\usepackage{fancyhdr}
\usepackage{titlesec}
\usepackage[hidelinks]{hyperref}
\usepackage[a4paper,margin=2.1cm,headheight=26pt,headsep=9pt]{geometry}

\definecolor{primary}{RGB}{150,98,162}
\definecolor{primarydark}{RGB}{92,52,110}
\definecolor{excol}{RGB}{0,135,90}
\definecolor{attncol}{RGB}{200,40,55}
\definecolor{methcol}{RGB}{90,90,100}

\newcommand{\runtitle}{Thème 1 — Difficile — Énoncés}
\pagestyle{fancy}\fancyhf{}
\fancyhead[L]{\small\textcolor{primarydark}{\textbf{Fiche 12} — Chimie organique}}
\fancyhead[R]{\small\textcolor{primarydark}{\runtitle}}
\fancyfoot[C]{\small\thepage}
\renewcommand{\headrulewidth}{0.8pt}
\renewcommand{\headrule}{\hbox to\headwidth{\color{primary}\leaders\hrule height \headrulewidth\hfill}}
\setlength{\parindent}{0pt}\setlength{\parskip}{3pt}

\tcbset{boxbase/.style={enhanced,breakable,boxrule=0.9pt,arc=2.5pt,
   left=3mm,right=3mm,top=2mm,bottom=2mm,fonttitle=\bfseries,coltitle=white,
   toptitle=0.5mm,bottomtitle=0.5mm}}
\newtcolorbox[auto counter]{exo}[1]{boxbase,colback=primary!4,colframe=primary,
   title={\textsc{Exercice}~\thetcbcounter\ \textmd{--- \itshape #1}}}
\newtcolorbox{consigne}{boxbase,colback=methcol!8,colframe=methcol,sharp corners,
   title={\textsc{Consignes}}}

\newcommand{\banner}[2]{%
\begin{tcolorbox}[enhanced,colback=primary,colframe=primary,arc=5pt,boxrule=0pt,
   left=5mm,right=5mm,top=2.5mm,bottom=2.5mm,coltext=white]
{\large\bfseries #1}\\[1pt]{\small #2}
\end{tcolorbox}\vspace{2pt}}

\begin{document}

\banner{Thème 1 — Hydrocarbures, formules brutes et squelette}{Niveau \textbf{Difficile} — Énoncés \quad$\vert$\quad Fiche 12, Chimie organique}

\begin{consigne}
Exercices \textbf{ouverts} à \textbf{sous-questions progressives} : chaque étape s'appuie sur la précédente. Formule du degré d'insaturation (hétéroatomes inclus) : $I=\dfrac{2n+2+p-m-q}{2}$ ($n$ : C ; $m$ : H ; $p$ : N ; $q$ : halogènes ; O sans effet).
\end{consigne}

% ---------------------------------------------------------------------
\begin{exo}{analyse complète d'un alcane ramifié}
On étudie le \textbf{2,2,4-triméthylpentane} (« isooctane », référence de l'indice d'octane des carburants) :
\begin{center}\chemfig{CH_3-[:30]C(-[:90]CH_3)(-[:150]CH_3)-[:-30]CH_2-[:30]CH(-[:90]CH_3)-[:-30]CH_3}\end{center}
\begin{enumerate}[label=\textbf{\alph*)},leftmargin=2em,itemsep=1pt]
  \item En complétant chaque carbone à 4 liaisons, détermine la \textbf{formule brute}.
  \item Calcule le degré d'insaturation $I$ et confirme qu'il s'agit d'un \textbf{alcane saturé}.
  \item Classe \emph{chacun} des carbones (primaire / secondaire / tertiaire / quaternaire) et donne le décompte total.
  \item Combien de carbones portent respectivement \textbf{3 H}, \textbf{2 H}, \textbf{1 H} et \textbf{0 H} ? Vérifie la cohérence avec la question c).
\end{enumerate}
\end{exo}

% ---------------------------------------------------------------------
\begin{exo}{hydrogénation et degré d'insaturation (d'après la colchicine)}
La \textbf{colchicine}, médicament de la crise de goutte, a pour formule brute $\ce{C22H25NO6}$. Après \textbf{hydrogénation complète} de toutes ses doubles liaisons $\ce{C=C}$ « réductibles », on obtient un composé de formule $\ce{C22H37NO6}$.
\begin{enumerate}[label=\textbf{\alph*)},leftmargin=2em,itemsep=1pt]
  \item Combien d'atomes d'hydrogène ont été ajoutés lors de l'hydrogénation ? En déduire le nombre de doubles liaisons $\ce{C=C}$ réduites (une $\ce{C=C}$ fixe une molécule $\ce{H2}$).
  \item Calcule le degré d'insaturation $I$ de la colchicine $\ce{C22H25NO6}$ (attention : l'azote intervient, l'oxygène non).
  \item Les insaturations \emph{non} réduites par l'hydrogénation ($I$ total $-$ nombre de $\ce{C=C}$ réduites) correspondent à des cycles, des carbonyles $\ce{C=O}$ et des liaisons de noyaux aromatiques. Combien en reste-t-il ?
\end{enumerate}
\end{exo}

% ---------------------------------------------------------------------
\begin{exo}{dénombrer et classer tous les isomères de \ce{C6H14}}
\begin{enumerate}[label=\textbf{\alph*)},leftmargin=2em,itemsep=1pt]
  \item Vérifie d'abord que $\ce{C6H14}$ est bien un \textbf{alcane} (calcul de $I$).
  \item Dessine (topologique ou semi-développée) et nomme \textbf{tous} les isomères de constitution de $\ce{C6H14}$. \emph{(Indication : il y en a 5.)}
  \item Pour chaque isomère, indique le décompte des carbones (I/II/III/IV). Lequel possède un carbone \textbf{quaternaire} ? Lequel possède \textbf{deux} carbones tertiaires ?
\end{enumerate}
\end{exo}

% ---------------------------------------------------------------------
\begin{exo}{degré d'insaturation avec hétéroatomes (paracétamol)}
Le \textbf{paracétamol} a pour formule brute $\ce{C8H9NO2}$. Sa molécule contient un \textbf{noyau benzénique}.
\begin{enumerate}[label=\textbf{\alph*)},leftmargin=2em,itemsep=1pt]
  \item Calcule son degré d'insaturation $I$ à l'aide de la formule générale (montre que le $\ce{O}$ n'entre pas dans le calcul et que le $\ce{N}$ y entre par $+p$).
  \item Un noyau benzénique compte pour combien d'insaturations ? Combien reste-t-il alors d'insaturations à « expliquer » ?
  \item Sachant que le paracétamol possède aussi une fonction \textbf{amide} ($\ce{-C(=O)-NH-}$), montre que le carbonyle de cette fonction rend compte de l'insaturation restante. Explique enfin \emph{pourquoi} un atome d'oxygène (divalent) n'ajoute jamais d'insaturation.
\end{enumerate}
\end{exo}

\end{document}
